%%%%%%%%%%%%%%%%%%%%% chapter-03.tex %%%%%%%%%%%%%%%%%%%%%%%%%%%%%%%%%
%
% Capítulo 3: Geometría computacional y formalización algebraica para la representación espacial
%
%%%%%%%%%%%%%%%%%%%%%%%%%%%%%%%%%%%%%%%%%%%%%%%%%%%%%%%%%%%%%%%%

\chapter{Geometría computacional y formación algebraica para la representación espacial}
\label{cap:geometria_computacional}

\abstract*{En este capítulo se establece la transición entre la abstracción lógica de los autómatas celulares y su representación geométrica tangible. Se formalizan las estructuras algebraicas que sustentan las retículas bidimensionales, definiendo el espacio discreto como un grupo abeliano bajo la adición. Asimismo, se analizan las métricas de distancia necesarias para delimitar vecindades y se desarrolla la geometría analítica de la tesela hexagonal. Este marco teórico constituye la base fundamental para el desarrollo del simulador en C/C++, permitiendo el mapeo riguroso de coordenadas discretas a espacios visuales continuos.}

\abstract{En este capítulo se establece la transición entre la abstracción lógica de los autómatas celulares y su representación geométrica tangible. Se formalizan las estructuras algebraicas que sustentan las retículas bidimensionales, definiendo el espacio discreto como un grupo abeliano bajo la adición. Asimismo, se analizan las métricas de distancia necesarias para delimitar vecindades y se desarrolla la geometría analítica de la tesela hexagonal. Este marco teórico constituye la base fundamental para el desarrollo del simulador en C/C++, permitiendo el mapeo riguroso de coordenadas discretas a espacios visuales continuos.}

\section{Estructuras algebraicas para retículas espaciales}
\label{sec:estructuras_algebraicas}

La comprensión de un autómata celular (AC) como sistema dinámico discreto requiere, en primera instancia, de una definición unívoca del espacio donde ocurre la evolución. Si bien en capítulos anteriores se abordó la dinámica de estados, en esta sección se formaliza el soporte topológico sobre el cual se asientan dichas celdas.

\subsection{Definición axiomática del espacio bidimensional discreto}

Se define la retícula celular, denotada por $\mathcal{L}$, como el conjunto de puntos en un espacio euclidiano que poseen coordenadas enteras. Para los fines de este Trabajo Terminal, se restringe el estudio al plano discreto bidimensional, el cual se identifica con el producto cartesiano de los enteros consigo mismo:

\begin{equation}
    \mathcal{L} = \mathbb{Z}^2 = \{ (i, j) \mid i, j \in \mathbb{Z} \}
\end{equation}

Cada par ordenado $c = (i, j) \in \mathbb{Z}^2$ representa la posición de una celda individual dentro del sistema. Es fundamental observar que este conjunto posee una estructura de \textit{grupo abeliano} bajo la operación de adición vectorial estándar. 

\begin{definition}
Sea $c_1 = (i_1, j_1)$ y $c_2 = (i_2, j_2)$ dos elementos de $\mathbb{Z}^2$. La operación suma $+ : \mathbb{Z}^2 \times \mathbb{Z}^2 \rightarrow \mathbb{Z}^2$ se define como:
\begin{equation}
    c_1 + c_2 = (i_1 + i_2, j_1 + j_2)
\end{equation}
\end{definition}

Esta estructura algebraica garantiza el cumplimiento de las siguientes propiedades fundamentales, las cuales son esenciales para la invarianza por traslación de las reglas de transición \cite{Kari2005}:

\begin{enumerate}
    \item \textbf{Clausura:} Para todo $c_1, c_2 \in \mathbb{Z}^2$, el resultado de $c_1 + c_2$ es un elemento que pertenece a $\mathbb{Z}^2$.
    \item \textbf{Asociatividad:} La agrupación de términos no altera el resultado de la suma de múltiples posiciones.
    \item \textbf{Existencia del elemento neutro:} Existe un vector origen $\mathbf{0} = (0, 0) \in \mathbb{Z}^2$ tal que para toda celda $c$, $c + \mathbf{0} = c$.
    \item \textbf{Existencia de elementos inversos:} Para cada celda $c = (i, j)$, existe un inverso $-c = (-i, -j)$ tal que $c + (-c) = \mathbf{0}$.
    \item \textbf{Conmutatividad:} El orden de la suma no altera la posición resultante, lo que simplifica el cálculo de desplazamientos relativos en la retícula.
\end{enumerate}

La naturaleza de grupo abeliano de $(\mathbb{Z}^2, +)$ permite que cualquier vecindad sea definida de manera relativa a una celda central mediante un conjunto fijo de vectores de desplazamiento, asegurando que la topología del espacio sea uniforme en todos sus puntos.

\subsection{Topología métrica de la red celular}

Para delimitar el alcance de la interacción entre celdas (vecindad), es imperativo introducir una función de distancia $d$ que actúe sobre la retícula. Una métrica en $\mathbb{Z}^2$ es una función $d: \mathbb{Z}^2 \times \mathbb{Z}^2 \rightarrow \mathbb{R}_{\geq 0}$ que satisface los axiomas de identidad de indiscernibles, simetría y desigualdad triangular.

En el contexto de los autómatas celulares, la elección de la métrica define la geometría de la vecindad. Se consideran tres métricas fundamentales inducidas por la norma $L_p$:

\begin{enumerate}
    \item \textbf{Distancia de Manhattan ($L_1$):} Define la distancia como la suma de las diferencias absolutas de sus coordenadas. Es la métrica natural para vecindades tipo Von Neumann.
    \begin{equation}
        d_1(c_1, c_2) = |i_1 - i_2| + |j_1 - j_2|
    \end{equation}
    
    \item \textbf{Distancia de Chebyshev ($L_\infty$):} Representa el máximo de las diferencias absolutas. Esta métrica es la base para la construcción de vecindades de Moore \cite{Wolfram2002}.
    \begin{equation}
        d_\infty(c_1, c_2) = \max(|i_1 - i_2|, |j_1 - j_2|)
    \end{equation}
    
    \item \textbf{Distancia Euclidiana ($L_2$):} Proporciona la distancia física en línea recta en el plano continuo $\mathbb{R}^2$.
    \begin{equation}
        d_2(c_1, c_2) = \sqrt{(i_1 - i_2)^2 + (j_1 - j_2)^2}
    \end{equation}
\end{enumerate}

La relevancia de estas métricas radica en la definición formal de la vecindad de radio $r$ para una celda $c$, denotada como $\mathcal{N}_r(c)$, la cual se establece como el conjunto de celdas cuya distancia a $c$ es menor o igual a $r$:

\begin{equation}
    \mathcal{N}_r(c) = \{ c' \in \mathbb{Z}^2 \mid d(c, c') \leq r \}
\end{equation}

Es posible observar que, para una retícula cuadrada convencional, una vecindad de Von Neumann corresponde a una bola cerrada bajo la métrica $L_1$, mientras que una vecindad de Moore corresponde a una bola cerrada bajo la métrica $L_\infty$. En el caso de la topología hexagonal que motiva este estudio, se requerirá una generalización de estas distancias para adaptarse a la simetría de seis ejes, lo cual se desarrollará analíticamente en la Sección 3.2.

% TODO: Insertar figura que compare gráficamente las vecindades de radio 1 bajo las métricas L1, L2 y Linf en una retícula cuadrada.
\begin{figure}[htbp]
    \centering
    % \includegraphics[width=0.7\textwidth]{metricas_vecindad}
    \caption{Representación geométrica de las vecindades inducidas por distintas métricas en $\mathbb{Z}^2$. Se observa cómo la elección de la distancia determina la morfología de la interacción local.}
    \label{fig:metricas_vecindad}
\end{figure}


\section{Geometría analítica de la tesela fundamental hexagonal}
\label{sec:geometria_analitica_hexagonal}

Tras la formalización de la retícula como un grupo abeliano, resulta imperativo caracterizar la geometría de la unidad mínima de información: la celda hexagonal. A diferencia de las teselas cuadrangulares, el hexágono regular ofrece una simetría superior y una equidistancia entre vecinos que favorece la emergencia de estructuras dinámicas complejas \cite{Basurto2009}.

\subsection{Parametrización geométrica de la celda unitaria en $\mathbb{R}^2$}

Sea $R$ el \textit{circunradio} de un hexágono regular, definido como la distancia desde el centro geométrico hacia cualquiera de sus vértices. Para propósitos de este modelado, se considera una orientación de ``punta superior'' (\textit{pointy-top}), donde el hexágono se encuentra centrado en el origen coordenado $(0,0)$.

Los seis vértices que delimitan la frontera de la celda, denotados por el conjunto $\mathcal{V} = \{V_0, V_1, \dots, V_5\}$, pueden ser calculados mediante identidades trigonométricas vectoriales. Dado que el hexágono regular divide el plano en ángulos de $\frac{\pi}{3}$ radianes ($60^\circ$), la posición de cada vértice $V_k$ está determinada por:

\begin{equation}
    V_k = \begin{pmatrix} x_k \\ y_k \end{pmatrix} = \begin{pmatrix} R \cos\left( \frac{\pi}{6} + \frac{k\pi}{3} \right) \\ R \sin\left( \frac{\pi}{6} + \frac{k\pi}{3} \right) \end{pmatrix}, \quad \text{para } k = 0, 1, \dots, 5.
\end{equation}

Es posible derivar la relación entre el circunradio $R$ y la \textit{apotema} (denotada como $a$), la cual representa la distancia mínima desde el centro hacia el punto medio de cualquiera de sus lados. Mediante el análisis del triángulo rectángulo formado por el radio, la apotema y la mitad de un lado, se establece que:

\begin{equation}
    a = R \cos\left( \frac{\pi}{6} \right) = \frac{\sqrt{3}}{2}R
\end{equation}

Esta constante geométrica es fundamental para el algoritmo de renderizado, pues define el factor de escala en el eje de las abscisas al realizar el empaquetamiento de la red.

% TODO: Insertar diagrama del hexágono regular centrado en el origen indicando R, a, y los ángulos de deflexión de los vértices.
\begin{figure}[htbp]
    \centering
    % \includegraphics[width=0.4\textwidth]{geometria_celda}
    \caption{Geometría analítica de la tesela hexagonal unitaria. Se destacan las relaciones trigonométricas entre el circunradio y la apotema.}
    \label{fig:geometria_celda}
\end{figure}

\subsection{Formulación del dominio poligonal y ecuaciones de contorno}

Para garantizar una detección rigurosa de pertenencia de puntos —proceso crítico en la interfaz gráfica del simulador—, el interior de la celda se modela como un conjunto convexo $\mathcal{H} \subset \mathbb{R}^2$. Este dominio se define matemáticamente como la intersección de seis semiespacios cerrados.

Sea $P = (x, y)$ un punto arbitrario en el plano. Dicho punto pertenece al dominio de la celda hexagonal si y solo si satisface el siguiente sistema de inecuaciones lineales simultáneas, derivadas de la simetría axial del polígono:

\begin{equation}
\label{eq:dominio_hexagonal}
\mathcal{H} = \left\{ (x, y) \in \mathbb{R}^2 \mid \begin{cases} 
|y| \leq \frac{\sqrt{3}}{2}R \\
\frac{\sqrt{3}}{2}|x| + \frac{1}{2}|y| \leq \frac{\sqrt{3}}{2}R 
\end{cases} \right\}
\end{equation}

La expresión \ref{eq:dominio_hexagonal} permite una verificación computacionalmente eficiente. La primera inecuación delimita la extensión vertical de la celda respecto a sus bordes horizontales, mientras que la segunda inecuación, de naturaleza escalar, restringe el espacio mediante las pendientes de las aristas laterales. 

Este modelo algebraico constituye la base analítica para el \textit{mapeo inverso} que se discutirá en secciones posteriores: dada una coordenada de pantalla producida por un evento de ratón, el sistema debe resolver este conjunto de inecuaciones para determinar la celda $(i, j)$ correspondiente en el espacio de estados. 

Resulta de interés observar que este planteamiento no depende de una búsqueda iterativa, sino de una evaluación directa de la norma del vector de posición, lo cual optimiza la complejidad temporal del simulador a tiempo constante $\mathcal{O}(1)$ por cada consulta de pertenencia.

% TODO: Insertar gráfica de las funciones de nivel que definen las inecuaciones del dominio hexagonal.
\begin{figure}[htbp]
    \centering
    % \includegraphics[width=0.6\textwidth]{interseccion_planos}
    \caption{Representación del dominio $\mathcal{H}$ como intersección de semiplanos. Las regiones sombreadas indican el cumplimiento de las restricciones lineales.}
    \label{fig:interseccion_planos}
\end{figure}



\section{Transformaciones afines y mapeo de coordenadas globales}
\label{sec:transformaciones_afines}

Una vez caracterizada la topología de la celda unitaria hexagonal, resulta imperativo establecer los mecanismos formales para proyectar una pluralidad de estas teselas sobre el plano euclidiano. La transición de la abstracción algebraica de la retícula $\mathbb{Z}^2$ hacia una representación gráfica continua en $\mathbb{R}^2$ requiere la definición de operadores matriciales que garanticen una biyección espacial sin superposiciones ni vacíos, logrando un teselado regular perfecto.

\subsection{Operadores matriciales para la traslación espacial}

Sea un par ordenado $(i,j) \in \mathbb{Z}^2$ que denota el índice discreto de una celda en el espacio de estados. Se define la transformación afín $\mathcal{T}: \mathbb{Z}^2 \to \mathbb{R}^2$ como la función que mapea las coordenadas lógicas hacia el centro geométrico continuo $(x,y)$ del hexágono correspondiente en el plano cartesiano.

Para una retícula hexagonal orientada en configuración de \textit{punta superior} (\textit{pointy-top}), los vectores base del espacio no son ortogonales. Un desplazamiento unitario en el eje discreto de las abscisas ($i$) corresponde a un movimiento completamente horizontal, mientras que un desplazamiento unitario en el eje de las ordenadas ($j$) implica un desplazamiento oblicuo.

Esta proyección se formaliza algebraicamente mediante la multiplicación de un vector de coordenadas lógicas por una matriz de transformación $M$, escalada por el circunradio $R$:

\begin{equation}
    \begin{pmatrix} x \\ y \end{pmatrix} = 
    R \begin{pmatrix} 
        \sqrt{3} & \frac{\sqrt{3}}{2} \\ 
        0 & \frac{3}{2} 
    \end{pmatrix} 
    \begin{pmatrix} i \\ j \end{pmatrix}
\end{equation}

Es posible observar que el determinante de la matriz de transformación es $\Delta = \frac{3\sqrt{3}}{2}$, un valor distinto de cero que garantiza la independencia lineal de los vectores base. Este modelo, conocido formalmente como \textit{sistema de coordenadas axiales}, es matemáticamente elegante debido a que preserva la estructura de grupo abeliano bajo la suma vectorial estándar. Sin embargo, su mapeo directo hacia arreglos de memoria bidimensionales en lenguajes como C/C++ genera un desperdicio de espacio debido al sesgo romboidal de la proyección \cite{Basurto2009}.

\subsection{Sistemas paramétricos de coordenadas compensadas (\textit{Offset Coordinates})}

Para resolver la disparidad entre la topología de las estructuras de datos ortogonales en la memoria computacional y la geometría no ortogonal de la retícula hexagonal, se introduce el sistema de \textit{coordenadas compensadas}. Este enfoque modela geométricamente el empaquetamiento óptimo de la malla aplicando una traslación condicionada por la paridad del índice de fila o columna.

Dado que la arquitectura topológica del simulador se orienta por filas horizontales, se define analíticamente el sistema compensado de fila impar (conocido en la literatura como sistema \textit{odd-r}). En este modelo, las celdas de las filas con índice $j$ impar sufren un desplazamiento horizontal equivalente a la apotema $a = \frac{\sqrt{3}}{2}R$.

Sea $P(j) = j \pmod 2$ la función indicadora de paridad booleana. El mapeo suprayectivo que determina las coordenadas euclidianas del centro del hexágono $(x, y)$ a partir de sus índices $(i,j)$ se define mediante el siguiente sistema de ecuaciones paramétricas:

\begin{equation}
    x(i, j) = \left( i + \frac{P(j)}{2} \right) \sqrt{3} R
\end{equation}

\begin{equation}
    y(i, j) = j \left( \frac{3}{2} R \right)
\end{equation}

Este sistema de ecuaciones revela las propiedades de empaquetamiento de la red. La distancia horizontal entre centros adyacentes en la misma fila es estrictamente $2a = \sqrt{3}R$, mientras que la distancia vertical entre filas consecutivas es $\frac{3}{2}R$, lo que provoca un anidamiento (\textit{interlocking}) de las celdas. Geométricamente, esta disposición es isomorfa al empaquetamiento denso de círculos en el plano bidimensional, asegurando que cada celda equidiste perfectamente de sus seis vecinos inmediatos \cite{Wolfram2002}.

Cabe destacar que, mediante rotaciones afines, es posible derivar el modelo complementario de columnas compensadas (\textit{even-q} u \textit{odd-q}) para hexágonos con orientación de borde plano (\textit{flat-top}), aunque para los fines de este Trabajo Terminal, la axiomatización se centra en la simetría horizontal dictada por el sistema \textit{odd-r}.

% TODO: Insertar un diagrama visual que muestre el mapeo de una matriz de memoria 2D ortogonal hacia una cuadrícula hexagonal bajo el sistema odd-r, destacando el desplazamiento (offset) en las filas impares.
\begin{figure}[htbp]
    \centering
    % \includegraphics[width=0.8\textwidth]{mapeo_compensado_oddr}
    \caption{Modelado geométrico del sistema de coordenadas compensadas \textit{odd-r}. Se ilustra cómo la paridad de la fila $j$ determina la traslación sobre el eje de las abscisas, emulando un empaquetamiento óptimo continuo a partir de índices discretos.}
    \label{fig:mapeo_compensado_oddr}
\end{figure}



\section{Topologías compactas y axiomatización de fronteras}
\label{sec:topologias_compactas}

La formulación teórica de un autómata celular asume, por definición, un espacio de estados infinito $\mathbb{Z}^2$. Sin embargo, cualquier emulación computacional está intrínsecamente restringida por una memoria finita. La delimitación abrupta de la retícula genera un conjunto de celdas fronterizas con vecindades incompletas, lo cual introduce perturbaciones (artefactos topológicos) que destruyen las estructuras dinámicas emergentes. Para preservar la simetría espacial y la computabilidad del modelo, resulta imperativo establecer una topología cerrada que elimine conceptualmente los límites del plano \cite{Toffoli1987}.

\subsection{Construcción algebraica del espacio cociente toroidal}

Sea una retícula finita de dimensiones $M \times N$, denotada por $\mathcal{L}_{M,N}$, donde los índices discretos operan en los rangos $i \in \{0, 1, \dots, M-1\}$ y $j \in \{0, 1, \dots, N-1\}$. Para estructurar la transición matemática de esta malla plana acotada hacia una variedad continua, se recurre a la teoría de clases de equivalencia mediante aritmética modular.

Se define una relación de equivalencia $\sim$ sobre el espacio infinito $\mathbb{Z}^2$ tal que dos coordenadas arbitrarias $(i, j)$ y $(i', j')$ son topológicamente indistinguibles si y solo si sus componentes son congruentes módulo las dimensiones de la retícula:

\begin{equation}
    (i, j) \sim (i', j') \iff i \equiv i' \pmod M \quad \land \quad j \equiv j' \pmod N
\end{equation}

Bajo esta relación, la clausura topológica de la matriz de estados finita se define como el espacio cociente:

\begin{equation}
    \mathbb{T}^2 \cong \mathbb{Z}^2 / (M\mathbb{Z} \times N\mathbb{Z})
\end{equation}

Este espacio cociente $\mathbb{T}^2$ es isomorfo a un toroide bidimensional discreto. Geométricamente, esta operación equivale a identificar el borde derecho de la matriz con el borde izquierdo (formando un cilindro), y posteriormente unir el borde superior con el inferior. La consecuencia axiomática fundamental de este mapeo es que toda celda en $\mathbb{T}^2$, sin excepción, posee una vecindad completa y simétrica. El espacio carece de fronteras, garantizando la invarianza por traslación de la función de transición global en la totalidad del dominio simulado.

% TODO: Insertar diagrama topológico que ilustre el pegado de las aristas opuestas de un plano rectangular para formar un toroide, señalando la continuidad de una vecindad hexagonal a través de la frontera (costura).
\begin{figure}[htbp]
    \centering
    % \includegraphics[width=0.6\textwidth]{espacio_toroidal}
    \caption{Representación del espacio cociente $\mathbb{T}^2$. Las fronteras opuestas se identifican algebraicamente para construir una topología continua libre de condiciones de borde.}
    \label{fig:espacio_toroidal}
\end{figure}

\subsection{Ecuaciones de transición para condiciones de contorno periódicas}

La viabilidad de un autómata celular complejo como sustrato de computación universal depende críticamente de la propagación ininterrumpida de información a través del espacio \cite{Wolfram2002}. Esta información viaja encapsulada en estructuras móviles (frecuentemente denominadas \textit{gliders}). Es indispensable demostrar que el mapeo toroidal conserva la cinemática de dichas partículas.

Sea $G(t) \subset \mathbb{T}^2$ el conjunto de celdas activas que conforman un glider en el instante $t$. Un glider se define rigurosamente por su periodo $p \in \mathbb{N}^+$ y su vector de traslación espacial $\vec{v} = (\Delta i, \Delta j)$, de tal forma que su evolución en un espacio infinito $\mathbb{Z}^2$ satisface $G(t+p) = G(t) + \vec{v}$ \cite{Basurto2009}.

Al confinar esta dinámica en el espacio cociente, la función de transición local $\delta$ evalúa las vecindades utilizando distancias geodésicas sobre el toroide. La posición de la estructura tras $p$ evoluciones, sujeta a las condiciones de contorno periódicas, está dictada por el mapeo biyectivo:

\begin{equation}
    G(t+p) = \left\{ \left( (i + \Delta i) \bmod M, \; (j + \Delta j) \bmod N \right) \mid (i, j) \in G(t) \right\}
\end{equation}

Dado que la operación módulo es distributiva respecto a la adición y preserva la distancia local $d(c_1, c_2)$ entre cualquier par de celdas vecinas a través de las "costuras" del toroide, la morfología de $G$ no experimenta distorsión topológica alguna al cruzar los límites $M$ o $N$. 

Matemáticamente, esto demuestra que el vector de traslación $\vec{v}$ se conserva de manera inercial. A largo plazo, cualquier glider que no colisione trazará una trayectoria cerrada cíclica (una curva geodésica discreta) sobre la variedad $\mathbb{T}^2$, retornando eventualmente a su configuración y posición inicial. Esta propiedad de \textit{ciclos límite} no solo valida la estabilidad del modelo computacional, sino que garantiza que las colisiones lógicas programadas entre múltiples gliders operen con determinismo estricto, independientemente de su localización global en el espacio simulado.



\section{Complejidad computacional del mapeo topológico y renderizado}
\label{sec:complejidad_renderizado}

La formalización algebraica y geométrica desarrollada en las secciones previas carecería de utilidad práctica si no pudiera ser proyectada y manipulada eficientemente en un entorno computacional. El desarrollo de un simulador visual en C/C++ requiere, por una parte, la capacidad de identificar estados a partir de la interacción continua del usuario (mapeo inverso) y, por otra, asegurar que el procesamiento paralelo inherente de los autómatas celulares se ejecute bajo cotas asintóticas viables. 

\subsection{Mapeo espacial inverso y localización de puntos}

Durante la ejecución del simulador, la interfaz gráfica permite la manipulación del espacio de estados mediante eventos de puntero. Esto plantea un problema de \textit{picking} geométrico: dado un punto continuo arbitrario $P = (x,y) \in \mathbb{R}^2$, es necesario determinar las coordenadas discretas $(i,j) \in \mathbb{Z}^2$ de la celda hexagonal que lo contiene. 

A diferencia de una retícula ortogonal donde la localización se resuelve mediante un truncamiento escalar directo, la teselación no ortogonal requiere la inversión del modelo afín. Partiendo del sistema de coordenadas axiales definido en la Sección 3.3, se establece la matriz de transformación continua-discreta mediante el cálculo de su inversa escalar:

\begin{equation}
    \begin{pmatrix} q_f \\ r_f \end{pmatrix} = 
    \frac{1}{R} \begin{pmatrix} 
        \frac{\sqrt{3}}{3} & -\frac{1}{3} \\ 
        0 & \frac{2}{3} 
    \end{pmatrix} 
    \begin{pmatrix} x \\ y \end{pmatrix}
\end{equation}

El resultado de esta operación proyecta el par bidimensional $(x,y)$ en un vector de coordenadas axiales fraccionarias $(q_f, r_f) \in \mathbb{R}^2$. Para obtener el índice discreto unívoco, se proyecta transitoriamente este punto bidimensional hacia un espacio cúbico tridimensional introduciendo una componente auxiliar $s_f = -q_f - r_f$. 

Al redondear matemáticamente $(q_f, r_f, s_f)$ al entero más cercano y ajustar la componente con mayor error de truncamiento para garantizar que $q + r + s = 0$, se obtiene la coordenada axial discreta. Finalmente, la restitución de esta terna hacia el sistema de memoria indexado por filas compensadas (\textit{odd-r}), denotado por $(i, j)$, se formaliza como:

\begin{equation}
    j = r, \quad i = q + \left\lfloor \frac{r - (r \bmod 2)}{2} \right\rfloor
\end{equation}

Este modelo algebraico inverso elimina por completo la necesidad de emplear algoritmos de búsqueda iterativa o iteraciones exhaustivas sobre el espacio de estados. La localización geométrica de cualquier perturbación externa se procesa en tiempo constante, sirviendo de fundamento matemático robusto para el motor gráfico en C/C++ \cite{Basurto2009}.

% TODO: Insertar diagrama conceptual del mapeo inverso. Se sugiere mostrar un vector (x,y) apuntando a una región interna de un hexágono y los pasos de transformación a través del plano axial hasta el índice (i,j).
\begin{figure}[htbp]
    \centering
    % \includegraphics[width=0.7\textwidth]{mapeo_inverso}
    \caption{Esquematización del mapeo algebraico inverso. La transición del dominio continuo de pantalla hacia el arreglo de memoria discreto.}
    \label{fig:mapeo_inverso}
\end{figure}

\subsection{Análisis espacial y temporal de la actualización topológica}

Para garantizar la viabilidad computacional de la simulación de colisiones y fenómenos complejos, es necesario someter los algoritmos de actualización global a un análisis de complejidad rigoroso. Como se establece en la teoría de análisis de algoritmos introducida por Cormen \textit{et al.} \cite{Cormen2009}, la notación asintótica proporciona una cota paramétrica sobre los recursos consumidos en relación al tamaño de la entrada.

Sea $n = M \times N$ la cardinalidad total de la retícula celular, representando el número de celdas en el espacio cerrado $\mathbb{T}^2$. La función de transición global requiere evaluar el vecindario de cada elemento. Dado que en una topología hexagonal el tamaño de la vecindad está estrictamente acotado a 6 vecinos colindantes (es decir, una constante independiente de $n$), el tiempo necesario para resolver las reglas lógicas locales es $\mathcal{O}(1)$. En consecuencia, la complejidad en tiempo para actualizar el estado del universo completo en una generación de simulación está dada por:

\begin{equation}
    T(n) = \sum_{k=1}^{n} \mathcal{O}(1) = \mathcal{O}(n)
\end{equation}

La cota superior lineal $\mathcal{O}(n)$ demuestra teóricamente que el algoritmo es asintóticamente óptimo para una arquitectura de procesamiento secuencial unihilo. 

Por otro lado, la semántica del autómata celular demanda un sincronismo perfecto; el estado futuro de la red $t+1$ debe evaluarse pura y exclusivamente sobre la configuración en el instante $t$. La actualización \textit{in-place} del arreglo de memoria violaría este principio, propagando perturbaciones espúreas. Por lo tanto, el sistema implementado en C/C++ requiere instanciar paralelamente un búfer secundario (técnica conocida como \textit{double buffering}). Esto establece que la complejidad espacial requerida por el motor matemático es:

\begin{equation}
    S(n) = \mathcal{O}(n)
\end{equation}

La inversión en requerimientos de memoria del modelo secuencial garantiza empíricamente la invariancia por traslación de la red y previene las condiciones de carrera lógicas, siendo el sustento informático indispensable para que la geometría analizada se traduzca correctamente en fenomenología computable.

%%%%%%%%%%%%%%%%%%%%% Fin del Capítulo 3 %%%%%%%%%%%%%%%%%%%%%%%%%%%%%%